\section{Related Work}

We situate our paper amongst existing work on vision-language compositionality, and debiasing datasets for model evaluation.

\textbf{Evaluating vision-language compositionality.}
Recent works have introduced benchmarks to evaluate the compositionality of vision-language models~\cite{RadfordKHRGASAM21}; they find that current models exhibit little compositional understanding~\cite{yuksekgonul2023when,thrush2022winoground,zhao2022vl, ma2022crepe, ray2023cola} despite their remarkable performance on downstream tasks~\cite{RadfordKHRGASAM21, 0001LXH22, singh2022flava, alayrac2022flamingo, wang2022omnivl,wang2022image, zhai2022lit}. 
Models have a hard time discerning between text containing the same words ordered differently~\cite{thrush2022winoground}. 
Models also fail to link objects to their attributes, or understand the relationship between objects~\cite{zhao2022vl, yuksekgonul2023when, ray2023cola}.
% Further, they perform even worse when compositions in the image are unseen during training or when the visual scenes and descriptions are more complex~\cite{ma2022crepe}.
Our work finds that many of the benchmarks used to evaluate compositionality have hackable biases; blind models that do not even look at the image outperform state-of-the-art vision-language models.

\textbf{Improving vision-language compositionality.} 
To enhance vision-language models’ compositionality, new proposals suggest training strategies that utilize additional data, models, and/or losses~\cite{yuksekgonul2023when, cascantebonilla2023going, ray2023cola, doveh2023teaching, Singh2023CoarsetoFineCL}. Amongst them, one prominent approach is to explicitly train the models to distinguish hard negatives from the correct captions~\cite{yuksekgonul2023when, doveh2023teaching}.
While these approaches appear to improve compositionality on benchmarks, it is unclear if these models achieve such improvements by actually acquiring compositional understanding or by exploiting biases in these datasets. We answer this question in our evaluation.


% \cy{Discuss that we focus on image-to-text. Add the problems that we observe.}


% \paragraph{Model-assisted dataset creation.} \cite{liu-etal-2022-wanli}

% \textbf{Improving vision-language compositionality.} 
% To improve the compositional understanding of VLMs, \cite{yuksekgonul2023when}, closer in spirit to ours, leverages hard negatives (which are composition-aware variants of the positive) in finetuning a CLIP model.
% Training with HNs~\cite{doveh2023teaching,yuksekgonul2023when}.
% Improving model compositionality~\cite{Singh2023CoarsetoFineCL}
% To enhance vision-language models’ compositionality, new proposals suggest training strategies that utilize additional data, models, and/or losses~\cite{yuksekgonul2023when, cascantebonilla2023going, doveh2023teaching, Singh2023CoarsetoFineCL}.
% For instance, NegCLIP~\cite{yuksekgonul2023when} adds hard negative captions and strong alternative images to the batch for loss calculation.
% SyViC~\cite{cascantebonilla2023going} finetunes vision-language models with synthetic image-text datasets built from a 3D physics-based simulation platform. MM-Adapter~\cite{ray2023cola} trains a lightweight multi-modal encoder on top of frozen disjoint image and text encoders with a finetuning image-text dataset containing descriptions of objects localized with attributes. 
% Others use an ``analogy loss'' between the ground truth caption and LLM-generated positives, and ``negative loss'' between the caption and generated negative texts~\cite{doveh2023teaching}.
% A coarse-to-fine contrastive loss has been introduced to incentivize alignment between an image and its descriptions of various complexities in addition to negative texts~\cite{Singh2023CoarsetoFineCL}. % (obtained from subgraphs of the complete scene graph)
% While these approaches appear to improve compositionality on benchmarks, it is unclear if these models achieve such improvements by actually acquiring compositional understanding or by exploiting biases in these datasets. We answer this question in our evaluation.

% \cy{Add that their performance is unclear on more robust dataset.} 


\textbf{Debiasing dataset for faithful model evaluation.}
Several prior manuscripts have pointed out that biased datasets could lead to an overestimation of models' true capabilities~\cite{gururangan-etal-2018-annotation}. They have proposed dataset de-biasing methods to enable more faithful model evaluations~\cite{reif2023fighting,zellers2018swag,sakaguchi2021winogrande,le2020adversarial}. For instance, adversarial filtering~\cite{zellers2018swag} iteratively trains an ensemble of classifiers on different training splits and uses them to filter out ``easy'' negatives for each instance.
% resulting in an ``adversarial'' dataset where the model suffers from high errors.
Building upon adversarial filtering, AFLite~\cite{sakaguchi2021winogrande,le2020adversarial} filters data instances in a more light-weight manner without retraining a model at each iteration and leads to benchmarks that more accurately represent the underlying tasks.
  % and is applied to create a challenging benchmark for commonsense reasoning. 
% Applying AFLITE filtering on existing vision and language benchmarks results in diverse examples that more accurately represent the underlying tasks~\cite{le2020adversarial}.
We use adversarial refinement to remove biases that creep into the generation of compositionality benchmarks.
 % Furthermore, a deeper dive into AFLITE reveals that it can be viewed as an approximation of the desirable but intractable optimum bias reduction objective, and applying AFLITE filtering on existing vision and language benchmarks results in diverse examples that more accurately represent the underlying tasks~\cite{le2020adversarial}.
% \cy{Fighting bias with bias \cite{reif2023fighting}}
